\documentclass[../../../include/open-logic-section]{subfiles}

\begin{document}

\olfileid[pt]{sth}{ord-arithmetic}{add}
\olsection{Adição Ordinal}

Suponha-se que se pretende somar $\alpha$ e $\beta$. Pode simplesmente colocar-se uma
{cópia} de $\beta$ imediatamente a seguir a uma cópia de $\alpha$. (É preciso tomar
\emph{cópias}, pois sabemos por
\olref[ordinals][basic]{ordinalsaresubsets} que ou $\alpha
\subseteq \beta$ ou $\beta \subseteq \alpha$.) O efeito intuitivo disto
é percorrer uma $\alpha$-sucessão de passos, \emph{e depois}
percorrer uma $\beta$-sucessão. A sucessão resultante fica
bem ordenada; logo, por
\olref[ordinals][ordtype]{thmOrdinalRepresentation} é isomorfa
a um ordinal (único). Esse ordinal pode considerar-se a \emph{soma}
de $\alpha$ e $\beta$.

Esta é a ideia intuitiva por detrás da adição ordinal. Para a definir
com rigor, começa-se pela ideia de tomar \emph{cópias} de conjuntos.
A ideia é usar etiquetas arbitrárias, $0$ e $1$, para saber de onde
veio cada objeto:

\begin{defn}\ollabel{defdissum}
A \emph{soma disjunta} de $A$ e $B$ é $A \disjointsum B = (A\times
\{0\}) \cup (B \times \{1\})$.
\end{defn}

Define-se em seguida uma ordem em pares de ordinais:

\begin{defn}
Para quaisquer ordinais $\alpha_1, \alpha_2, \beta_1, \beta_2$, diga-se que:
\begin{align*}
	\tuple{\alpha_1, \alpha_2} \rlexless \tuple{\beta_1, \beta_2}\text{ sse }&
	\text{ou $\alpha_2 \in \beta_2$}\\
	& \text{ou tanto $\alpha_2 = \beta_2$ como $\alpha_1 \in \beta_1$}
\end{align*}
\end{defn}

Trata-se de uma ordem \emph{lexicográfica inversa}, pois ordena-se pelo
segundo elemento e depois pelo primeiro. Recorde-se que se pretendia
definir $\alpha \ordplus \beta$ como o tipo de ordem de uma cópia de $\alpha$
seguida de uma cópia de $\beta$. Para o conseguir, define-se:

\begin{defn}\ollabel{defordplus}
Para quaisquer ordinais $\alpha$, $\beta$, a sua soma é $\alpha \ordplus
\beta = \ordtype{\alpha \disjointsum \beta, \rlexless}$.
\end{defn}
\noindent
Note-se que aqui se abusou ligeiramente da notação; estritamente deveria escrever-se ``$\Setabs{\tuple{x,y}\in \alpha\disjointsum\beta}{x \rlexless y}$'' em lugar de ``$\rlexless$''. Por brevidade, contudo, continuaremos a abusar da notação deste modo no que se segue.

O resultado seguinte, juntamente com
\olref[ordinals][ordtype]{thmOrdinalRepresentation}, confirma que a nossa
definição está bem formada:

\begin{lem}\ollabel{ordsumlessiswo}
$\tuple{\alpha \disjointsum \beta, \rlexless}$ é um bom ordenamento, para
quaisquer ordinais $\alpha$ e $\beta$.
\end{lem}

\begin{proof}
Obviamente $\rlexless$ é conexa em $\alpha \disjointsum \beta$. Para
mostrar que é bem fundada, fixe-se um $X \subseteq \alpha
\disjointsum \beta$ não vazio. Seja $Y$ o subconjunto de $X$ cuja segunda coordenada é a menor possível, isto é, $Y = \Setabs{\tuple{\gamma, i} \in X}{(\forall \tuple{\delta, j} \in X)i \leq j}$. Escolha-se agora o elemento de $Y$ com menor primeira coordenada.
\end{proof}
\noindent
Temos assim uma definição explícita e clara da adição ordinal. Eis um
fato não surpreendente (recorde-se que  $1 = \{0\}$, por
\olref[z][infinity-again]{defnomega}):
\begin{prop}
$\alpha \ordplus  1 = \ordsucc{\alpha}$, para qualquer ordinal $\alpha$.
\end{prop}

\begin{proof}
Considere-se o isomorfismo $f$ de $\ordsucc{\alpha} = \alpha \cup
\{\alpha\}$ para $\alpha\disjointsum1 = (\alpha \times \{0\})
\disjointsum (\{0\} \times \{1\})$ dado por $f(\gamma) =
\tuple{\gamma, 0}$ para $\gamma \in \alpha$, e $f(\alpha) = \tuple{0,
1}$.
\end{proof}
\noindent
Além disso, é fácil mostrar que a adição obedece a certas condições
recursivas:

\begin{lem}\ollabel{ordadditionrecursion}
Para quaisquer ordinais $\alpha, \beta$, tem-se:
\begin{align*}
	\alpha\ordplus 0 &= \alpha\\
	\alpha \ordplus  (\beta\ordplus 1) &= (\alpha \ordplus  \beta) \ordplus  1\\
	\alpha  \ordplus  \beta &= \supstrict_{\delta < \beta}(\alpha \ordplus  \delta) && \text{se $\beta $ é um ordinal limite}
\end{align*}
\end{lem}

\begin{proof}
Verifica-se caso a caso; primeiro:
\begin{align*}
	\alpha \ordplus  0
	& = \ordtype{(\alpha \times \{0\}) \cup (0 \times \{1\}), \rlexless} \\
	&= \ordtype{(\alpha \times \{0\}) \cup \{0\}, \rlexless}\\
	&= \alpha\\
	\alpha \ordplus (\beta \ordplus  1)
	%&= \ordtype{(\alpha\times \{0\}) \cup (\ordtype{\beta \ordplus  1}\times \{1\}), \rlexless} \\
	&= \ordtype{(\alpha\times \{0\}) \cup (\ordsucc{\beta}\times \{1\}), \rlexless} \\
	&= \ordtype{(\alpha\times \{0\}) \cup (\beta \times \{1\}), \rlexless} \ordplus 1\\
	&= (\alpha \ordplus  \beta) \ordplus  1
\end{align*}
Seja agora $\beta \neq \emptyset$ um limite. Se $\delta < \beta$ então
também $\delta\ordplus 1 < \beta$, logo $\alpha \ordplus  \delta$ é um
segmento inicial próprio de $\alpha \ordplus  \beta$. Logo $\alpha
\ordplus  \beta$ é um majorante estrito de $X = \Setabs{\alpha
\ordplus  \delta}{\delta < \beta}$. Além disso, se $\alpha \leq \gamma <
\alpha \ordplus  \beta$, então claramente $\gamma = \alpha \ordplus
\delta$ para algum $\delta < \beta$. Logo $\alpha \ordplus  \beta =
\supstrict_{\delta< \beta}(\alpha\ordplus \delta)$.
\end{proof}

Mas eis um fato marcante. Para definir a adição ordinal, poder-se-ia
\emph{em alternativa} ter usado simplesmente o teorema da recursão transfinita, e
estipulado as equações recursivas, exatamente como em
\olref{ordadditionrecursion} (embora com ``$\ordsucc{\beta}$''
em vez de ``$\beta \ordplus 1$'').

Há, pois, duas maneiras diferentes de definir operações nos
ordinais. Podem definir-se \emph{sinteticamente}, construindo explicitamente um conjunto bem ordenado e considerando o seu tipo de ordem. Ou
podem definir-se \emph{recursivamente}, apenas estipulando as equações
recursivas. Feito corretamente, porém, o resultado é idêntico. De fato,
\olref[ordinals][ordtype]{thmOrdinalRepresentation} garante
que estas equações recursivas determinam ordinais \emph{únicos}.

Em muitos aspectos, a aritmética ordinal comporta-se como a adição dos
números naturais. Por exemplo, pode provar-se o seguinte:

\begin{lem}\ollabel{ordinaladditionisnice}
Se $\alpha, \beta, \gamma$ são ordinais, então:
\begin{enumerate}
	\item\ollabel{ordaddition1} se $\beta < \gamma$, então $\alpha
	\ordplus  \beta < \alpha \ordplus  \gamma$
	\item\ollabel{ordaddition2} se $\alpha \ordplus  \beta =
	\alpha\ordplus \gamma$, então $\beta = \gamma$
	\item\ollabel{ordaddition3}  $\alpha \ordplus  (\beta \ordplus
	\gamma) = (\alpha \ordplus  \beta) \ordplus  \gamma$, isto é,
	a adição é associativa
	\item\ollabel{ordaddition4}  Se $\alpha \leq \beta$, então $\alpha
	\ordplus  \gamma \leq \beta \ordplus \gamma$
\end{enumerate}
\end{lem}

\begin{proof}
Provamos \olref{ordaddition3}, deixando o resto como exercício. A
prova é por indução transfinita simples em $\gamma$, usando
\olref{ordadditionrecursion}. Quando $\gamma = 0$:
\[
(\alpha \ordplus  \beta) \ordplus  0 = \alpha \ordplus  \beta  = \alpha \ordplus  (\beta \ordplus  0)
\]
Quando $\gamma = \delta\ordplus 1$, suponha-se por indução que $(\alpha
\ordplus  \beta) \ordplus  \delta = \alpha \ordplus  (\beta \ordplus
\delta)$; usando \olref{ordadditionrecursion} três vezes:
\begin{align*}
	(\alpha \ordplus  \beta) \ordplus  (\delta \ordplus  1) & = ((\alpha \ordplus  \beta) \ordplus  \delta)\ordplus 1\\
	& = (\alpha \ordplus  (\beta \ordplus  \delta)) \ordplus  1\\
	& = \alpha \ordplus  ((\beta \ordplus  \delta)\ordplus 1)\\
	& = \alpha \ordplus  (\beta \ordplus  (\delta\ordplus 1))
\end{align*}
Quando $\gamma$ é um ordinal limite, suponha-se por indução que se
$\delta \in \gamma$ então $(\alpha \ordplus  \beta) \ordplus  \delta =
\alpha \ordplus  (\beta \ordplus  \delta)$; então:
\begin{align*}
	(\alpha \ordplus  \beta) \ordplus  \gamma & = \supstrict_{\delta < \gamma}((\alpha \ordplus  \beta) \ordplus  \delta) \\
	&= \supstrict_{\delta < \gamma}(\alpha \ordplus  (\beta \ordplus  \delta))\\
	&= \alpha \ordplus  \supstrict_{\delta < \gamma}(\beta \ordplus  \delta)\\
	& = \alpha \ordplus  (\beta \ordplus  \gamma)
\end{align*}
\end{proof}

\begin{prob}
Prove o resto de
\olref[sth][ord-arithmetic][add]{ordinaladditionisnice}.
\end{prob}

Deste modo, a adição ordinal deve parecer muito familiar. Há contudo um ponto crucial em que a adição ordinal \emph{não} se parece com a adição nos naturais.

\begin{prop}\ollabel{ordsumnotcommute}
A adição ordinal \emph{não} é comutativa; $1 \ordplus  \omega = \omega <
\omega \ordplus  1$.
\end{prop}

\begin{proof}
Note-se que $1 \ordplus  \omega = \supstrict_{n < \omega} (1 \ordplus n)
= \omega \in \omega \cup \{\omega\} = \ordsucc{\omega} = \omega
\ordplus  1$.
\end{proof}
\noindent
Embora isto possa surpreender à primeira vista, \emph{não devia}. Por
um lado, ao considerar $1 \ordplus  \omega$, pensa-se no
tipo de ordem que se obtém colocando um elemento extra \emph{antes} de
todos os números naturais. Raciocinando como com o Hotel de Hilbert em
\olref[sfr][infinite][hilbert]{sec}, intuitivamente esse primeiro elemento
extra não altera o tipo de ordem global. Por
outro lado, ao considerar $\omega \ordplus  1$, pensa-se no
tipo de ordem que se obtém colocando um elemento extra \emph{depois} de
todos os números naturais. E isso é algo radicalmente diferente!

\end{document}
